\chapter*{Preface}
\addcontentsline{toc}{chapter}{Preface}
\markboth{Preface}{Preface}

The Gujarat State Eligibility Test (GSET) is the gateway to the post of Assistant Professor in the universities and colleges of Gujarat. Every candidate, whatever the subject, has to clear Paper~I, the common paper on teaching and research aptitude. Year after year, aspirants discover that the surest way to prepare for Paper~I is to work through the previous papers: the same concepts, the same types of reasoning puzzles and even the same questions return again and again.

This book brings together \textbf{twenty complete GSET Paper~I question papers}, from January 2002 to November 2025, with \textbf{1{,}070 questions} in all. For every question you will find:
\begin{itemize}
\item the full question exactly as set, with \textbf{all four options};
\item the \textbf{correct answer} according to the \emph{final} answer key published by the GSET Agency;
\item a \textbf{detailed explanation} of why the answer is correct and, wherever it helps, why the other options are wrong;
\item \textbf{redrawn diagrams}: Venn diagrams, number figures, dice, pie charts, bar graphs, line graphs and data tables, all typeset afresh so that they are clear in print;
\item \textbf{memory tips} that connect the question to the wider syllabus.
\end{itemize}

The official keys are followed faithfully. On the few occasions where an official key is clearly inconsistent with its own question, the book still records the official answer but adds a short \emph{Note on the key} explaining the correct reasoning, so that you learn the concept correctly and also know what the Agency accepted. Questions that the Agency dropped (marked Z or X) are clearly flagged and still solved.

The papers are arranged from the most recent to the oldest, because the recent papers reflect the current syllabus and pattern (including the NEP~2020 era). The older papers remain valuable: the same concepts in teaching, research, communication, reasoning and higher education appear throughout.

I hope this book gives you both the knowledge and the confidence to clear GSET. Best wishes for your examination and your teaching career.

\bigskip
\hfill\textit{Yash Kavaiya}\\
\hspace*{\fill}Gujarat, 2026
